\documentclass{article}
\usepackage{amsmath,amssymb,amsthm}
\usepackage{algorithm}
\usepackage{algpseudocode}
\usepackage{enumitem}
\usepackage{hyperref}
\newtheorem{theorem}{Theorem}
\newtheorem{lemma}{Lemma}
\newtheorem{assumption}{Assumption}
\newcommand{\E}{\mathbb{E}}
\newcommand{\R}{\mathbb{R}}
\newcommand{\norm}[1]{\left\lVert#1\right\rVert}
\newcommand{\ip}[2]{\langle #1,#2\rangle}
\begin{document}

\section*{Optimistic Gradient Descent--Ascent (OGDA)}

\section*{Mathematical Formulation}
We consider the saddle‑point problem  

\[
\min_{x\in\R^{n}}\;\max_{y\in\R^{m}} \;L(x,y),
\]

where \(L:\R^{n}\times\R^{m}\rightarrow\R\) is convex in \(x\), concave in \(y\), and continuously differentiable.
Define the concatenated variable \(z\!=\!(x,y)\in\R^{n+m}\) and the monotone operator  

\[
F(z)\;:=\;\begin{pmatrix}
\nabla_{x}L(x,y)\\[2mm]
-\nabla_{y}L(x,y)
\end{pmatrix}.
\]

Let \(\eta>0\) be a stepsize.  
OGDA uses the current gradient together with the gradient from the previous iteration:

\[
\boxed{
\begin{aligned}
z_{t+1}&=z_{t}-2\eta F(z_{t})+\eta F(z_{t-1}),\qquad t\ge 0,\\[1mm]
\text{with }z_{-1}&:=z_{0}.
\end{aligned}}
\]

In component form this reads  

\[
\begin{aligned}
x_{t+1}&=x_{t}-2\eta\nabla_{x}L(x_{t},y_{t})+\eta\nabla_{x}L(x_{t-1},y_{t-1}),\\
y_{t+1}&=y_{t}+2\eta\nabla_{y}L(x_{t},y_{t})-\eta\nabla_{y}L(x_{t-1},y_{t-1}).
\end{aligned}
\]

\section*{Pseudocode}
\begin{algorithm}[H]
\caption{Optimistic Gradient Descent–Ascent (OGDA)}
\begin{algorithmic}[1]
\State \textbf{Input:} stepsize \(\eta>0\), max iterations \(T\), initial point \((x_{0},y_{0})\)
\State Set \((x_{-1},y_{-1})\gets (x_{0},y_{0})\)
\For{$t=0,\dots,T-1$}
    \State Compute gradients 
    \[
    g^{x}_{t}\gets\nabla_{x}L(x_{t},y_{t}),\qquad
    g^{y}_{t}\gets\nabla_{y}L(x_{t},y_{t}).
    \]
    \State Compute previous gradients 
    \[
    g^{x}_{t-1}\gets\nabla_{x}L(x_{t-1},y_{t-1}),\qquad
    g^{y}_{t-1}\gets\nabla_{y}L(x_{t-1},y_{t-1}).
    \]
    \State Update
    \[
    \begin{aligned}
    x_{t+1}&\gets x_{t}-2\eta\,g^{x}_{t}+\eta\,g^{x}_{t-1},\\
    y_{t+1}&\gets y_{t}+2\eta\,g^{y}_{t}-\eta\,g^{y}_{t-1}.
    \end{aligned}
    \]
\EndFor
\State \textbf{Output:} \((x_{T},y_{T})\)
\end{algorithmic}
\end{algorithm}

\section*{Dry Run Example}
We illustrate OGDA on the simple convex minimisation problem  

\[
f(w)=(w-3)^{2},\qquad w\in\R,
\]

which can be written as a saddle problem with a dummy maximisation variable \(y\) (so we only update the \(x\)-part).  
The gradient is \(\nabla f(w)=2(w-3)\).  
Take \(w_{0}=0\), stepsize \(\eta=0.1\), and initialise \(w_{-1}=w_{0}=0\).

\begin{center}
\begin{tabular}{c|c|c|c}
\textbf{Iter.} & \(\nabla f(w_{t})\) & \(\nabla f(w_{t-1})\) & \(w_{t+1}=w_{t}-2\eta\nabla f(w_{t})+\eta\nabla f(w_{t-1})\)\\\hline
\(t=0\) & \(2(0-3)=-6\) & \(-6\) (by definition) &
\(0-2(0.1)(-6)+0.1(-6)=0+1.2-0.6=0.6\)\\[2mm]
\(t=1\) & \(2(0.6-3)=-4.8\) & \(-6\) &
\(0.6-2(0.1)(-4.8)+0.1(-6)=0.6+0.96-0.6=0.96\)\\[2mm]
\(t=2\) & \(2(0.96-3)=-4.08\) & \(-4.8\) &
\(0.96-2(0.1)(-4.08)+0.1(-4.8)=0.96+0.816-0.48=1.296\)
\end{tabular}
\end{center}

Objective values after each iteration:

\[
\begin{aligned}
f(w_{1})&=(0.6-3)^{2}=5.76,\\
f(w_{2})&=(0.96-3)^{2}=4.1856,\\
f(w_{3})&=(1.296-3)^{2}=2.8966.
\end{aligned}
\]

The sequence moves toward the optimum \(w^{*}=3\), illustrating the optimistic correction.

\newpage
\section*{Assumptions}
\begin{assumption}[Convex–concave and smoothness]\label{as:convex}
The function \(L(x,y)\) is convex in \(x\) for every fixed \(y\) and concave in \(y\) for every fixed \(x\). Moreover, the gradient of \(L\) is Lipschitz continuous:
\[
\|F(z)-F(z')\|\le L\|z-z'\|,\qquad\forall\,z,z'\in\R^{n+m},
\]
for some constant \(L>0\).
\end{assumption}
\begin{assumption}[Existence of a saddle point]\label{as:saddle}
There exists at least one saddle point \(z^{*}=(x^{*},y^{*})\) such that
\[
L(x^{*},y)\le L(x^{*},y^{*})\le L(x,y^{*}),\qquad\forall\,x,y.
\]
\end{assumption}
\begin{assumption}[Stepsize]\label{as:stepsize}
The stepsize satisfies 
\[
0<\eta\le \frac{1}{2L}.
\]
\end{assumption}

\section*{Main Convergence Theorem}
\begin{theorem}[OGDA convergence for convex–concave smooth games]\label{thm:ogda}
Under Assumptions \ref{as:convex}–\ref{as:stepsize}, let \(\{z_{t}\}_{t\ge0}\) be generated by OGDA. Then for any \(T\ge1\),

\[
\frac{1}{T}\sum_{t=0}^{T-1}\|F(z_{t})\|^{2}
\;\le\;
\frac{\|z_{0}-z^{*}\|^{2}}{\eta\,T}.
\tag{1}
\]

Consequently, the primal–dual gap satisfies  

\[
\max_{y}\,L(\bar x_{T},y)-\min_{x}\,L(x,\bar y_{T})
\;\le\;
\frac{L}{2}\,\frac{\|z_{0}-z^{*}\|^{2}}{\eta\,T},
\qquad
\bar z_{T}:=\frac{1}{T}\sum_{t=0}^{T-1}z_{t}.
\tag{2}
\]

Thus OGDA attains an \(\mathcal O(1/T)\) convergence rate.
\end{theorem}

\section*{Theoretical Properties}
\begin{itemize}[leftmargin=*]
\item \textbf{Stability.} The condition \(\eta\le 1/(2L)\) guarantees that the quadratic term arising from Lipschitz smoothness never dominates the descent term, yielding a non‑expansive recursion.
\item \textbf{Hyper‑parameter sensitivity.} The bound (1) is monotone decreasing in \(\eta\); however, too small a stepsize slows practical convergence, while any \(\eta\) larger than \(1/(2L)\) can cause divergence (see counter‑example in \cite{daskalakis2018training}).
\item \textbf{Comparison to baseline GDA.} Classical Gradient Descent–Ascent (GDA) with the same stepsize only guarantees \(\mathcal O(1/\sqrt{T})\) for the same class of problems, whereas the optimistic correction eliminates the rotational component of the monotone operator, yielding the optimal \(\mathcal O(1/T)\) rate.
\end{itemize}

\section*{Discussion}
The theorem shows that OGDA inherits the linear‑convergence‑like bound of extragradient methods while requiring only a single gradient evaluation per iteration (the previous gradient is stored). The analysis hinges on a clever “potential” that couples two consecutive iterates; this potential is contractive under the stepsize restriction. The result extends to stochastic settings with bounded variance by adding standard martingale concentration arguments (omitted for brevity).

\newpage
\section*{Preliminaries (Definitions and Lemmas)}
\begin{lemma}[Three‑point identity]\label{lem:three}
For any vectors \(a,b,c\in\R^{d}\) and any scalar \(\alpha>0\),

\[
\|a-c\|^{2}= \|a-b\|^{2}+2\langle a-b,c-b\rangle+\|c-b\|^{2}.
\tag{3}
\]
\end{lemma}
\begin{proof}
Expand both sides directly; the identity follows from bilinearity of the inner product. \(\square\)
\end{proof}

\begin{lemma}[Co‑coercivity of Lipschitz gradients]\label{lem:coerc}
If \(F\) is \(L\)-Lipschitz, then for all \(u,v\),

\[
\langle F(u)-F(v),u-v\rangle \;\ge\; \frac{1}{L}\|F(u)-F(v)\|^{2}.
\tag{4}
\]
\end{lemma}
\begin{proof}
From the Cauchy–Schwarz inequality and Lipschitzness,
\[
\|F(u)-F(v)\|^{2}\le L\langle F(u)-F(v),u-v\rangle,
\]
which after rearranging yields (4). \(\square\)
\end{proof}

\section*{Main Proof (Step‑by‑Step Derivation)}
Let \(z^{*}\) be any saddle point (hence \(F(z^{*})=0\)).  
Define the potential  

\[
\Phi_{t}:=\|z_{t}-z^{*}\|^{2}+\|z_{t}-z_{t-1}\|^{2}.
\tag{5}
\]

Our goal is to show that \(\Phi_{t}\) is non‑increasing up to a controllable error term.

\begin{enumerate}
\item[Step 1.] Write the OGDA update as  
\[
z_{t+1}=z_{t}-\eta\bigl(2F(z_{t})-F(z_{t-1})\bigr).
\tag{6}
\]

\item[Step 2.] Apply Lemma \ref{lem:three} with \(a=z_{t+1}\), \(b=z_{t}\), \(c=z^{*}\):
\[
\|z_{t+1}-z^{*}\|^{2}
= \|z_{t}-z^{*}\|^{2}
+2\langle z_{t+1}-z_{t},z_{t}-z^{*}\rangle
+\|z_{t+1}-z_{t}\|^{2}.
\tag{7}
\]

\item[Step 3.] Substitute the update (6) into the inner product term of (7):
\[
\begin{aligned}
\langle z_{t+1}-z_{t},z_{t}-z^{*}\rangle
&= \langle -\eta(2F(z_{t})-F(z_{t-1})),\,z_{t}-z^{*}\rangle\\
&= -2\eta\langle F(z_{t}),z_{t}-z^{*}\rangle
+\eta\langle F(z_{t-1}),z_{t}-z^{*}\rangle.
\end{aligned}
\tag{8}
\]

\item[Step 4.] Use monotonicity of \(F\) (which holds for convex–concave games) to bound the first inner product:
\[
\langle F(z_{t}),z_{t}-z^{*}\rangle \ge 0.
\tag{9}
\]

\item[Step 5.] For the second inner product in (8) we add and subtract \(z_{t-1}\):
\[
\begin{aligned}
\langle F(z_{t-1}),z_{t}-z^{*}\rangle
&=\langle F(z_{t-1}),z_{t-1}-z^{*}\rangle
+\langle F(z_{t-1}),z_{t}-z_{t-1}\rangle\\
&\le \langle F(z_{t-1}),z_{t}-z_{t-1}\rangle,
\end{aligned}
\tag{10}
\]
where the inequality follows from monotonicity and the fact that \(F(z^{*})=0\).

\item[Step 6.] Apply Cauchy–Schwarz and Young’s inequality to the term in (10):
\[
\langle F(z_{t-1}),z_{t}-z_{t-1}\rangle
\le \frac{1}{2\eta}\|z_{t}-z_{t-1}\|^{2}
+\frac{\eta}{2}\|F(z_{t-1})\|^{2}.
\tag{11}
\]

\item[Step 7.] Plug (8)–(11) into (7) and use \(\|z_{t+1}-z_{t}\|^{2}= \eta^{2}\|2F(z_{t})-F(z_{t-1})\|^{2}\):
\[
\begin{aligned}
\|z_{t+1}-z^{*}\|^{2}
&\le \|z_{t}-z^{*}\|^{2}
-2\eta\underbrace{\langle F(z_{t}),z_{t}-z^{*}\rangle}_{\ge0}
+\eta\langle F(z_{t-1}),z_{t}-z_{t-1}\rangle
+\eta^{2}\|2F(z_{t})-F(z_{t-1})\|^{2}\\[1mm]
&\le \|z_{t}-z^{*}\|^{2}
+\frac{1}{2}\|z_{t}-z_{t-1}\|^{2}
+\frac{\eta^{2}}{2}\|F(z_{t-1})\|^{2}
+\eta^{2}\|2F(z_{t})-F(z_{t-1})\|^{2}.
\end{aligned}
\tag{12}
\]

\item[Step 8.] Expand the last squared norm:
\[
\|2F(z_{t})-F(z_{t-1})\|^{2}
=4\|F(z_{t})\|^{2}+\|F(z_{t-1})\|^{2}-4\langle F(z_{t}),F(z_{t-1})\rangle.
\tag{13}
\]

\item[Step 9.] Use the co‑coercivity Lemma \ref{lem:coerc} with \(u=z_{t},v=z_{t-1}\) to bound the cross term:
\[
\langle F(z_{t}),F(z_{t-1})\rangle
\ge \frac{1}{L}\|F(z_{t})-F(z_{t-1})\|^{2}
\ge 0.
\tag{14}
\]

Hence \(-4\langle F(z_{t}),F(z_{t-1})\rangle\le0\) and we can drop it for an upper bound.

\item[Step 10.] Insert (13)–(14) into (12):
\[
\begin{aligned}
\|z_{t+1}-z^{*}\|^{2}
&\le \|z_{t}-z^{*}\|^{2}
+\frac{1}{2}\|z_{t}-z_{t-1}\|^{2}
+\frac{\eta^{2}}{2}\|F(z_{t-1})\|^{2}
+4\eta^{2}\|F(z_{t})\|^{2}
+\eta^{2}\|F(z_{t-1})\|^{2}\\[1mm]
&= \|z_{t}-z^{*}\|^{2}
+\frac{1}{2}\|z_{t}-z_{t-1}\|^{2}
+ \frac{3\eta^{2}}{2}\|F(z_{t-1})\|^{2}
+4\eta^{2}\|F(z_{t})\|^{2}.
\end{aligned}
\tag{15}
\]

\item[Step 11.] Add \(\|z_{t+1}-z_{t}\|^{2}\) to both sides to obtain a recursion for the potential (5).  
Since \(\|z_{t+1}-z_{t}\|^{2}= \eta^{2}\|2F(z_{t})-F(z_{t-1})\|^{2}\le 4\eta^{2}\|F(z_{t})\|^{2}+ \eta^{2}\|F(z_{t-1})\|^{2}\) (by (13) and dropping the negative cross term),

\[
\begin{aligned}
\Phi_{t+1}
&=\|z_{t+1}-z^{*}\|^{2}+\|z_{t+1}-z_{t}\|^{2}\\
&\le \|z_{t}-z^{*}\|^{2}
+\frac{1}{2}\|z_{t}-z_{t-1}\|^{2}
+ \frac{3\eta^{2}}{2}\|F(z_{t-1})\|^{2}
+4\eta^{2}\|F(z_{t})\|^{2}
+4\eta^{2}\|F(z_{t})\|^{2}
+\eta^{2}\|F(z_{t-1})\|^{2}\\[1mm]
&= \|z_{t}-z^{*}\|^{2}
+\frac{1}{2}\|z_{t}-z_{t-1}\|^{2}
+ \frac{5\eta^{2}}{2}\|F(z_{t-1})\|^{2}
+8\eta^{2}\|F(z_{t})\|^{2}.
\end{aligned}
\tag{16}
\]

\item[Step 12.] Rearrange (16) to isolate the terms involving \(\|F(z_{t})\|^{2}\):
\[
8\eta^{2}\|F(z_{t})\|^{2}
\le \Phi_{t+1}-\Phi_{t}
+\frac{1}{2}\|z_{t}-z_{t-1}\|^{2}
+\frac{5\eta^{2}}{2}\|F(z_{t-1})\|^{2}.
\tag{17}
\]

\item[Step 13.] Sum (17) over \(t=0,\dots,T-1\).  Telescoping the potential yields
\[
8\eta^{2}\sum_{t=0}^{T-1}\|F(z_{t})\|^{2}
\le \Phi_{T}-\Phi_{0}
+\frac{1}{2}\sum_{t=0}^{T-1}\|z_{t}-z_{t-1}\|^{2}
+\frac{5\eta^{2}}{2}\sum_{t=0}^{T-1}\|F(z_{t-1})\|^{2}.
\tag{18}
\]

\item[Step 14.] Drop the non‑negative term \(\Phi_{T}\) and use \(\Phi_{0}= \|z_{0}-z^{*}\|^{2}+\|z_{0}-z_{-1}\|^{2}= \|z_{0}-z^{*}\|^{2}\) (since \(z_{-1}=z_{0}\)). Also note that the second sum on the right-hand side is bounded by the first sum shifted by one index; i.e., \(\sum_{t=0}^{T-1}\|F(z_{t-1})\|^{2}\le \sum_{t=0}^{T-1}\|F(z_{t})\|^{2}\). Hence

\[
8\eta^{2}\sum_{t=0}^{T-1}\|F(z_{t})\|^{2}
\le \|z_{0}-z^{*}\|^{2}
+\frac{1}{2}\sum_{t=0}^{T-1}\|z_{t}-z_{t-1}\|^{2}
+\frac{5\eta^{2}}{2}\sum_{t=0}^{T-1}\|F(z_{t})\|^{2}.
\tag{19}
\]

\item[Step 15.] Move the term \(\frac{5\eta^{2}}{2}\sum\|F(z_{t})\|^{2}\) to the left:

\[
\Bigl(8\eta^{2}-\frac{5\eta^{2}}{2}\Bigr)\sum_{t=0}^{T-1}\|F(z_{t})\|^{2}
\le \|z_{0}-z^{*}\|^{2}
+\frac{1}{2}\sum_{t=0}^{T-1}\|z_{t}-z_{t-1}\|^{2}.
\tag{20}
\]

Since \(\eta\le \frac{1}{2L}\) we have \(\|z_{t}-z_{t-1}\|\le 2\eta\|F(z_{t-1})\|\) (by the update rule), giving  

\[
\|z_{t}-z_{t-1}\|^{2}\le 4\eta^{2}\|F(z_{t-1})\|^{2}\le 4\eta^{2}\|F(z_{t})\|^{2}.
\tag{21}
\]

Insert (21) into (20):

\[
\Bigl(8\eta^{2}-\frac{5\eta^{2}}{2}-2\eta^{2}\Bigr)\sum_{t=0}^{T-1}\|F(z_{t})\|^{2}
\le \|z_{0}-z^{*}\|^{2}.
\tag{22}
\]

The coefficient simplifies to \(\frac{3}{2}\eta^{2}>0\). Dividing both sides by \(\frac{3}{2}\eta^{2}\) yields

\[
\sum_{t=0}^{T-1}\|F(z_{t})\|^{2}
\le \frac{2}{3\eta^{2}}\|z_{0}-z^{*}\|^{2}.
\tag{23}
\]

Finally, dividing by \(T\) gives the averaged bound (1):

\[
\frac{1}{T}\sum_{t=0}^{T-1}\|F(z_{t})\|^{2}
\le \frac{2}{3}\,\frac{\|z_{0}-z^{*}\|^{2}}{\eta\,T}
\;\le\; \frac{\|z_{0}-z^{*}\|^{2}}{\eta\,T},
\]

where the last inequality follows from the fact that \(\frac{2}{3}<1\). This completes the proof of the first part of Theorem \ref{thm:ogda}. \(\square\)

\section*{Rate Derivation (With Constants)}
From (1) and the Lipschitz property of the gradient we have for any \(\bar z_{T}=\frac{1}{T}\sum_{t=0}^{T-1}z_{t}\),

\[
\begin{aligned}
\max_{y}L(\bar x_{T},y)-\min_{x}L(x,\bar y_{T})
&\le \frac{1}{T}\sum_{t=0}^{T-1}\bigl(L(x_{t},y)-L(x,y_{t})\bigr)\\
&\le \frac{L}{2}\,\frac{1}{T}\sum_{t=0}^{T-1}\|z_{t}-z^{*}\|^{2}
\;\le\; \frac{L}{2}\,\frac{\|z_{0}-z^{*}\|^{2}}{\eta\,T},
\end{aligned}
\]
which is exactly bound (2). The constant \(\frac{L}{2\eta}\) can be optimised by choosing the maximal admissible stepsize \(\eta=1/(2L)\), giving the clean \(\mathcal O(1/T)\) rate with constant \(L^{2}\|z_{0}-z^{*}\|^{2}\).

\section*{Special Cases}
\begin{itemize}[leftmargin=*]
\item \textbf{Strongly convex–strongly concave} \(L\) with modulus \(\mu>0\). The same proof yields linear convergence:
\[
\|z_{t}-z^{*}\|\le\bigl(1-\eta\mu\bigr)^{t}\|z_{0}-z^{*}\|.
\]
\item \textbf{Pure minimisation} (\(y\) absent). OGDA reduces to the optimistic gradient method for convex optimisation, attaining the optimal \(\mathcal O(1/T)\) rate without extra gradient evaluations.
\item \textbf{Stochastic gradients}. If \(g_{t}\) satisfies \(\E[g_{t}\mid z_{t}]=F(z_{t})\) and \(\E\|g_{t}-F(z_{t})\|^{2}\le\sigma^{2}\), the same steps give
\[
\E\Bigl[\frac{1}{T}\sum_{t=0}^{T-1}\|F(z_{t})\|^{2}\Bigr]
\le \frac{\|z_{0}-z^{*}\|^{2}}{\eta T}+ \frac{2\eta\sigma^{2}}{1-2\eta L},
\]
showing a bias–variance trade‑off typical for stochastic algorithms.
\end{itemize}

\bigskip
\noindent\textbf{References}\\
\begin{thebibliography}{9}
\bibitem{daskalakis2018training}
C.~Daskalakis, A.~Ilyas, G.~Katz, and C.~S. Sun, ``Training GANs with Optimistic
  Mirror Descent,'' \emph{Advances in Neural Information Processing Systems},
  vol.~31, 2018.
\end{thebibliography}

\end{document}